% !TeX program = lualatex
% Compile twice: lualatex 77.tex
% Standalone research notebook; original section preserved.
% Research additions: 27 September 2026. No external TeX input or BibTeX file.
\documentclass[11pt,a4paper]{article}
\usepackage[margin=24mm,headheight=14pt]{geometry}
\usepackage{fontspec}
\setmainfont{Latin Modern Roman}
\setsansfont{Latin Modern Sans}
\setmonofont{Latin Modern Mono}
\usepackage{amsmath,amssymb,mathtools}
\usepackage{unicode-math}
\setmathfont{Latin Modern Math}
\usepackage{microtype}
\usepackage{enumitem,xurl,needspace}
\usepackage[dvipsnames]{xcolor}
\usepackage{hyperref}
\hypersetup{unicode=true,colorlinks=true,linkcolor=MidnightBlue,urlcolor=MidnightBlue,
 pdftitle={Research notebook - Problem 77},pdfauthor={Open Problems Math},bookmarksnumbered=true}
\usepackage{bookmark}
\usepackage{titlesec}
\titleformat{\section}{\Large\bfseries\raggedright}{\thesection}{0.7em}{}
\usepackage{fancyhdr}
\pagestyle{fancy}\fancyhf{}
\fancyhead[L]{\small\sffamily Open Problems Math | Research notebook}
\fancyhead[R]{\small\sffamily Problem 77 | 27 September 2026}
\fancyfoot[C]{\thepage}
\setlength{\parindent}{0pt}
\setlength{\parskip}{5pt plus 1pt}
\setlength{\emergencystretch}{3em}
\setcounter{tocdepth}{1}
\setcounter{secnumdepth}{1}
\setlist[itemize]{leftmargin=3em,itemsep=5pt,topsep=4pt,parsep=0pt}
\allowdisplaybreaks
\numberwithin{equation}{section}
\newcommand{\N}{\mathbb{N}}
\newcommand{\Z}{\mathbb{Z}}
\newcommand{\Q}{\mathbb{Q}}
\newcommand{\R}{\mathbb{R}}
\newcommand{\C}{\mathbb{C}}
\newcommand{\F}{\mathbb{F}}
\newcommand{\A}{\mathbb{A}}
\newcommand{\PP}{\mathbb P}
\newcommand{\E}{\mathbb{E}}
\newcommand{\eps}{\varepsilon}
\newcommand{\dd}{\,\mathrm d}
\newcommand{\sref}[2]{\hyperlink{source:#1:#2}{[S#2]}}
\newcommand{\eref}[2]{\hyperlink{extra:#1:#2}{[E#2]}}
\newif\ifshowcatalogue
\showcataloguetrue
\newcommand{\cataloguescope}[1]{\ifshowcatalogue
 \paragraph{Exact scope in the supplied catalogue.}
 \begin{quote}\small #1\end{quote}\fi}
\begin{document}
\setcounter{section}{76}
\section{\texorpdfstring{Chowla conjecture on correlations of the Liouville function}{Chowla conjecture on correlations of the Liouville function}}\label{77}
\subsection{Definitions and mathematical statement}
For $n=\prod_p p^{v_p(n)}$, put $\Omega(n)=\sum_pv_p(n)$ and $\lambda(n)=(-1)^{\Omega(n)}$. For the positive-slope linear forms $L_i(n)=a_in+b_i$ specified in the record, require pairwise nonproportionality $a_ib_j-a_jb_i\ne0$. The assertion is
\[
 \frac1x\sum_{1\le n\le x}\prod_{i=1}^k\lambda(a_in+b_i)\longrightarrow0.
\]
This is an ordinary, unweighted Cesàro average. Replacing it by a logarithmically weighted average changes the claim. The record's distinct nonnegative offsets and positive natural slopes are retained.
\par\smallskip\textit{Formulation and concept references:} \sref{77}{1}.
\cataloguescope{Prove that for the Liouville function lambda, for every k >= 2, distinct nonnegative integers b\_1,...,b\_k and natural numbers a\_1,...,a\_k with a\_i b\_j - a\_j b\_i != 0 for i<j, sum\_\{n<=x\} lambda(a\_1 n + b\_1) ... lambda(a\_k n + b\_k) = o(x) as x -> infinity.}
\subsection{Short English statement}
Do the parities of prime-factor counts decorrelate along every allowed finite family of distinct linear patterns?
% BEGIN RESEARCH_ADDITION ATTEMPT_77
\subsection{Research attempt}

\paragraph{Current frontier.}
Tao's two-point theorem is logarithmically averaged \eref{77}{1}; the
logarithmic equivalences in \sref{77}{1} do not convert it into the
ordinary average required here. Chinis's 2026 result assumes Siegel zeros
and obtains correlations at infinitely many scales, not all large scales
unconditionally \eref{77}{2}. The dynamical averages in \sref{77}{4} also
retain their averaging qualifications. The attempted upgrade below is a
Tauberian one: identify the extra control that would exclude rare
macroscopic failures of ordinary cancellation.

\paragraph{Attack route.}
Partial summation relates ordinary and logarithmic means, but the inverse
implication can fail. A second route would control an average square of
ordinary prefix sums on logarithmic scales. The latter is stronger than
what the target demands, but has a useful feature: bounded summands prevent
a large normalized prefix sum from occupying arbitrarily narrow scale
intervals. That persistence can turn an integrable energy bound into
pointwise cancellation. A counterexample to the simpler logarithmic route
is proved first, rather than concealed.

\paragraph{Attempt: an exact separation of averaging modes.}
For any bounded complex sequence $a_n$ put
$A(x)=\sum_{n\leq x}a_n$ and $L(x)=\sum_{n\leq x}a_n/n$.
Partial summation gives, for $x\geq1$,
\begin{equation}\label{eq77:abel}
 L(x)=\frac{A(x)}x+\int_1^x\frac{A(t)}{t^2}\,dt.
\end{equation}
Now take $a_n=(-1)^{\lfloor\log_2n\rfloor}$.
At the ends of complete dyadic blocks,
\begin{equation}\label{eq77:counter}
 A(2^K-1)=\sum_{j=0}^{K-1}(-2)^j
           =\frac{1-(-2)^K}{3}.
\end{equation}
The ratios $A(2^K-1)/(2^K-1)$ have limiting values $1/3$ and $-1/3$
along alternate parities. Nevertheless $L(x)=O(1)$.
Indeed the harmonic sum on block $[2^j,2^{j+1})$ is
$\log2+O(2^{-j})$; the alternating main terms are bounded and the errors
are absolutely summable. A final incomplete block contributes at most
$\log2+O(1)$. Thus even bounded logarithmic partial sums, much stronger
than $L(x)=o(\log x)$, need not force the ordinary mean to vanish.
This is an artificial sequence, not a multiplicative-function
counterexample, and precisely rejects an unqualified Tauberian inference.

\paragraph{Attempt: a finite-energy sufficient lemma.}
Assume $|a_n|\leq1$ and
\begin{equation}\label{eq77:energy}
 \int_2^\infty\frac{|A(t)|^2}{t^3}\,dt<\infty.
\end{equation}
Then $A(x)=o(x)$. To prove this, suppose instead that
$|A(x)|\geq\delta x$ for some fixed $\delta>0$ and arbitrarily large $x$.
We may assume $0<\delta\leq1$. On
$[x,x+\delta x/4]$, bounded increments imply
\[
 |A(t)-A(x)|\leq t-x+1\leq\delta x/4+1.
\]
For sufficiently large $x$ this gives $|A(t)|\geq\delta x/2$ throughout
that interval. Since $t\leq5x/4$, its contribution to the integral in
\eqref{eq77:energy} is at least
\[
 \frac{\delta x}{4}\,
 \frac{\delta^2x^2/4}{(5x/4)^3}
 =\frac{4\delta^3}{125}.
\]
From an unbounded sequence of such $x$, choose successively disjoint
intervals. Their positive contributions contradict finiteness of the
integral. This proves the lemma, including the possible complex phase of
$A(x)$: only the triangle inequality was used.

For the original fixed tuple set
$a_n=\prod_{i=1}^k\lambda(a_i n+b_i)$; the slopes and offsets guarantee
positive arguments, and $|a_n|=1$. Thus establishing
\eqref{eq77:energy} for every allowed tuple would prove exactly its ordinary
Chowla correlation assertion. No uniformity in the tuple is required by
this sufficient condition; the integral and its bound may depend on all
the fixed coefficients. The method still asks for stronger quantitative
summability than $A(x)=o(x)$ itself, so it is not asserted to be equivalent
to Chowla.

\paragraph{Attempt: making the missing correlation estimate explicit.}
For a finite real $Y>2$, expansion of the finite sum gives
\begin{equation}\label{eq77:kernel}
 \int_2^Y\frac{|A(t)|^2}{t^3}\,dt
 =\sum_{m,n\leq Y}a_m\overline{a_n}
 \left(\frac{1}{2\max(2,m,n)^2}-\frac{1}{2Y^2}\right).
\end{equation}
For each pair the kernel is the integral from $\max(2,m,n)$ to $Y$;
terms at the upper endpoint have zero contribution. Finite summation
justifies the identity without any rearrangement of a conditionally
convergent infinite double series.

The desired arithmetic input is boundedness of the right side as
$Y\to\infty$, using cancellation between $m$ and $n$.
The diagonal terms are summable, but estimating all terms in absolute
value gives only an $O(\log Y)$ bound. Indeed there are $O(j)$ pairs with
$\max(m,n)=j$, each of weight $O(j^{-2})$.
For the constant sequence $a_n=1$, that logarithmic divergence actually
occurs, since $A(t)\asymp t$. Therefore positivity of the quadratic kernel
and bounded summands alone do not remove it.

One possible continuation is a bilinear estimate on separated blocks of
$m,n$ that is summable over scales. It must apply to the product of
Liouville values on the original linear patterns, not to independently
resampled values. Replacing the off-diagonal terms by zero because
``Liouville is random'' assumes the intended cancellation. Applying a
one-variable multiplicative mean-value theorem is also unjustified: a
shifted product $\prod_i\lambda(a_i n+b_i)$ is generally not a
multiplicative function of $n$.

\paragraph{Stress test.}
The energy condition excludes repeated large spikes by paying a fixed
amount for each spike, but it may fail for sequences whose normalized
sums tend slowly to zero. Consequently failure of this sufficient condition
would not disprove Chowla. A logarithmic average controls a signed first
moment on scales; it may cancel between large positive and negative blocks,
as \eqref{eq77:counter} shows. It is not an estimate for the nonnegative
energy in \eqref{eq77:energy}.

A bound established only on some scales under an extra zero hypothesis
cannot be applied as a bound for every finite $Y$ in
\eqref{eq77:kernel}. Nor can a statement averaged over coefficients or
offsets supply a conclusion for each fixed tuple. The companion script
verifies \eqref{eq77:counter} with exact integers and the finite kernel
identity with rational arithmetic on small sequences; it does not test
the infinite energy of any unresolved Liouville pattern.

\paragraph{Outcome.}
\textbf{Outcome: Conditional reduction.}

A fully proved finite-energy Tauberian lemma would upgrade suitable
quantitative correlation control to the ordinary target, while an explicit
bounded sequence disproves the simpler logarithmic upgrade. The exact
finite kernel identifies the missing two-variable cancellation and avoids
an invalid infinite-sum interchange.

The energy estimate is not proved for the required Liouville products,
so no new Chowla case or counterexample is claimed. A complete continuation
must establish summable scale control or a weaker Tauberian condition
actually available arithmetically. The elementary lemmas carry no claim
of independent novelty.

% Keep the local bibliography together rather than leave a source fragment.
\needspace{170mm}
% END RESEARCH_ADDITION ATTEMPT_77
\subsection{Sources}
\begin{itemize}\raggedright
\item[\textnormal{[S1]}]\hypertarget{source:77:1}{} Terence Tao, Equivalence of the logarithmically averaged Chowla and Sarnak conjectures, abstract (2016).
\url{https://arxiv.org/abs/1605.04628}.
{\small\textit{Catalogue formulation source.}}
\item[\textnormal{[S2]}]\hypertarget{source:77:2}{} The Chowla conjecture and Landau–Siegel zeroes.
\url{https://www.cambridge.org/core/journals/mathematical-proceedings-of-the-cambridge-philosophical-society/article/chowla-conjecture-and-landausiegel-zeroes/515F3378450DED201A21E26BF801CEEB}.
{\small\textit{Catalogue research/status reference; not a proof endorsement.}}
\item[\textnormal{[S3]}]\hypertarget{source:77:3}{} From Bateman-Horn to Chowla.
\url{https://www.preprints.org/manuscript/202601.0277}.
{\small\textit{Catalogue research/status reference; not a proof endorsement.}}
\item[\textnormal{[S4]}]\hypertarget{source:77:4}{} Two averaged dynamical generalizations of Chowla's conjecture.
\url{https://arxiv.org/abs/2608.16108}.
{\small\textit{Catalogue research/status reference; not a proof endorsement.}}
% BEGIN RESEARCH_ADDITION REFERENCES_77
\item[\textnormal{[E1]}]\hypertarget{extra:77:1}{} Terence Tao,
\emph{The logarithmically averaged Chowla and Elliott conjectures for
two-point correlations}, Forum of Mathematics, Pi 4 (2016), e8;
arXiv:1509.05422, main theorem and logarithmic averaging convention.
\url{https://arxiv.org/abs/1509.05422}.
{\small\textit{Inspected 27 September 2026; not an ordinary Ces\`aro theorem.}}
\item[\textnormal{[E2]}]\hypertarget{extra:77:2}{} Jake Chinis,
\emph{Siegel zeros and Sarnak's conjecture}, Journal of the London
Mathematical Society 113 (2026), no. 3, e70479.
\url{https://londmathsoc.onlinelibrary.wiley.com/doi/10.1112/jlms.70479}.
{\small\textit{Inspected 27 September 2026; the Siegel-zero hypothesis and
infinitely-many-scales conclusion are both retained.}}
% END RESEARCH_ADDITION REFERENCES_77
\end{itemize}

\end{document}
